\VerdierPhysicalPageAnchor{0059}%
\VerdierSemanticAnchor{ch1-1-5-13}
\noindent\textbf{1.5.13.}\quad
The $\mathbb Z^n$-categories that we shall use in this work will be
strict categories. We shall therefore not attempt to address the problem
indicated in (1.5.12) in the general case. However, the example presented
in (1.5.12) suggests that determining the structures of a
$\mathbb Z$-category on a category $\mathcal C$ is particularly simple.
This is indeed the case. Let $E\text{-}\operatorname{Cat}$ denote the
$U$-category whose objects are the categories belonging to $U$ equipped
with an equivalence:

\VerdierSemanticAnchor{ch1-1-5-13-eq1}
\begin{equation}
  T\colon\mathcal C\longrightarrow\mathcal C
  \tag{1.5.13.1}
\end{equation}
and whose morphisms are the functors
$F\colon\mathcal C\to\mathcal C'$ equipped with an isomorphism:

\VerdierSemanticAnchor{ch1-1-5-13-eq2}
\begin{equation}
  m\colon FT\longrightarrow T'F.
  \tag{1.5.13.2}
\end{equation}

There is then evidently a functor from the $U$-category
$\mathbb Z\text{-}\operatorname{Cat}$ to the $U$-category
$E\text{-}\operatorname{Cat}$:

\VerdierSemanticAnchor{ch1-1-5-13-eq3}
\begin{equation}
  R\colon \mathbb Z\text{-}\operatorname{Cat}
  \longrightarrow E\text{-}\operatorname{Cat}.
  \tag{1.5.13.3}
\end{equation}

\VerdierSemanticAnchor{ch1-1-5-13-action}
The functor $R$ assigns to a $\mathbb Z$-category
$(\mathcal C,\mathbb Z,T,c)$ the category $\mathcal C$ equipped with
the equivalence $T(1)\colon\mathcal C\to\mathcal C$, and to every
$\mathbb Z$-functor
\[
  (F,m)\colon(\mathcal C,\mathbb Z,T,c)
  \longrightarrow(\mathcal C',\mathbb Z,T',c'),
\]
the functor $F$ equipped with the isomorphism
$m(1)\colon FT(1)\to T'(1)F$.

\VerdierSemanticAnchor{ch1-1-5-14}
\VerdierSemanticAnchor{ch1-1-5-14-proposition}
\noindent\textbf{Proposition 1.5.14.}\quad
\emph{The functor $R$ is an equivalence of categories.}

\VerdierSemanticAnchor{ch1-1-5-14-proof}
We shall confine ourselves to giving indications for the proof.
Proposition (1.5.2) and Remark (1.5.8) show that $R$ is fully faithful.
It therefore suffices to show that $R$ is essentially surjective. Thus,
let $(\mathcal C,T\colon\mathcal C\to\mathcal C)$ be a category
equipped with an equivalence. Let
$T^{-1}\colon\mathcal C\to\mathcal C$ be a quasi-inverse functor to
$T$, and let

\VerdierSemanticAnchor{ch1-1-5-14-eq1}
\begin{equation}
  \begin{aligned}
    u&\colon\operatorname{id}_{\mathcal C}\longrightarrow TT^{-1},\\
    v&\colon\operatorname{id}_{\mathcal C}\longrightarrow T^{-1}T
  \end{aligned}
  \tag{1.5.14.1}
\end{equation}
be the adjunction isomorphisms.

For every integer $n>0$, the functors $T^n$ and $(T^{-1})^n$ are
adjoint to one another, whence the adjunction isomorphisms

\VerdierSemanticAnchor{ch1-1-5-14-eq2}
\begin{equation}
  \begin{aligned}
    u^{(n)}&\colon\operatorname{id}_{\mathcal C}
      \longrightarrow T^n(T^{-1})^n,\\
    v^{(n)}&\colon\operatorname{id}_{\mathcal C}
      \longrightarrow (T^{-1})^nT^n.
  \end{aligned}
  \tag{1.5.14.2}
\end{equation}
